\documentclass[10pt]{article}

\usepackage{hyperref}
\usepackage{xcolor}
\usepackage{calc}
\usepackage{graphicx}
\usepackage{tikz}
\usepackage{fontspec}
\usepackage{fontawesome5}
\usepackage{titlesec}
\usepackage{enumitem}
\usepackage{fancybox}
\usepackage{xeCJK}
\usepackage[
    a4paper,
    left=1.1cm,
    right=1.1cm,
    top=0.78cm,
    bottom=0.88cm,
    nohead
]{geometry}

\hypersetup{hidelinks}
\urlstyle{same}
\setlength{\parindent}{0pt}
\setlength{\parskip}{0pt}
\pagenumbering{gobble}
\linespread{1.08}
\renewcommand{\arraystretch}{1.12}

\IfFontExistsTF{Noto Serif CJK SC}{
    \setCJKmainfont{Noto Serif CJK SC}
}{
    \IfFontExistsTF{Microsoft YaHei}{
        \setCJKmainfont{Microsoft YaHei}
    }{
        \setCJKmainfont{SimSun}
    }
}

\definecolor{primarycolor}{RGB}{200,60,60}
\definecolor{secondarycolor}{RGB}{200,68,28}
\definecolor{headercolor}{RGB}{152,44,42}

\setlist[itemize]{
    leftmargin=1.35em,
    itemsep=0.12em,
    topsep=0.18em,
    parsep=0pt,
    partopsep=0pt
}
\setlist[enumerate]{leftmargin=*}

\newlength{\iconwidth}
\setlength{\iconwidth}{1.5em}

\titleformat{\section}
  {\Large\bfseries\raggedright}
  {}{0em}
  {}
  [{\color{secondarycolor}\titlerule}]
\titlespacing*{\section}{0pt}{0.55em}{0.28em}

\newcommand{\sectionicon}[2]{\makebox[\iconwidth][c]{\color{primarycolor}{#1}}\quad #2}
\newcommand{\resumeitem}[2]{\item \textbf{#1}：#2}
\newcommand{\projectheading}[3]{%
    \noindent\makebox[\textwidth][s]{%
        \makebox[0pt][l]{{\large \textbf{#1}}}%
        \hfill
        \makebox[0pt][c]{{\normalsize \textbf{#2}}}%
        \hfill
        \makebox[0pt][r]{#3}%
    }%
}

\newcommand{\contact}{
    \footnotesize
    \textcolor{white}{
        \faEnvelope \quad 邮箱：\href{mailto:email@example.com}{email@example.com}
        \hspace{2.4em}
        \faPhone \quad 电话：13800000000
        \hspace{2.4em}
        \faGithub \quad 主页：\href{https://github.com/example}{github.com/example}
    }
}

\begin{document}

    \begin{tikzpicture}[remember picture, overlay]
        \fill[headercolor] (current page.north west) rectangle ([yshift=-1.05cm]current page.north east);
        \node[anchor=north] at ([yshift=-0.24cm]current page.north){\contact};
    \end{tikzpicture}
    \vspace{-1.0em}

    \begin{minipage}[t]{0.79\textwidth}
        \section{\sectionicon{\faAddressCard}{个人信息}}
        \begin{tabular*}{\textwidth}{@{\extracolsep{\fill}}lll}
            \textbf{姓\qquad 名}：候选人姓名 & \textbf{出生年月}：2000.01 & \textbf{籍\qquad 贯}：城市 \\
            \textbf{学\qquad 校}：学校名称 & \textbf{学\qquad 历}：硕士/本科 & \textbf{求职方向}：AI 应用开发
        \end{tabular*}
        \vspace{0.45em}

        \section{\sectionicon{\faGraduationCap}{教育背景}}
        {\textbf{学校名称}}，专业名称，\textit{硕士/本科} \hfill 2024.09--2027.06
        \begin{itemize}
            \resumeitem{荣誉奖项}{奖学金、竞赛、排名或其他可信成果}
        \end{itemize}
    \end{minipage}
    \hfill
    \begin{minipage}[t]{0.14\textwidth}
        \vspace{-0.35em}
        \setlength{\fboxsep}{0pt}
        % 如果没有照片，删除整个右侧 minipage，并把左侧 minipage 调整为 1.0\textwidth。
        \IfFileExists{images/avatar.jpg}{
            \doublebox{\includegraphics[width=\linewidth]{images/avatar.jpg}}
        }{
            \doublebox{\parbox[c][2.6cm][c]{\linewidth}{\centering 照片}}
        }
    \end{minipage}
    \vspace{0.45em}

    \section{\sectionicon{\faWrench}{专业技能}}
    \begin{itemize}
        \resumeitem{AI 应用开发}{熟悉 RAG、工具调用、多 Agent 协作与结构化输出流程，能够围绕业务场景设计 LLM 应用链路。}
        \resumeitem{模型与推理优化}{了解 Transformer、Prompt Engineering、SFT/LoRA 等方法，具备模型 API 集成或推理部署经验。}
        \resumeitem{工程能力}{掌握 Python、数据库、Linux、Docker、Git、接口联调与性能测试，具备完整工程交付能力。}
        \resumeitem{全栈协作}{熟悉前后端分离开发，并能使用 AI 编程工具提升需求分析、开发、测试和交付效率。}
    \end{itemize}

    \vspace{0.45em}

    \section{\sectionicon{\faBriefcase}{实习经历}}
    \projectheading{公司名称}{岗位名称}{2025.06--2025.09}
    \begin{itemize}
        \resumeitem{核心职责}{覆盖需求分析、数据处理、模型调研、Prompt/RAG 流程优化、接口设计、测试、服务封装及部署联调等环节。}
        \resumeitem{关键成果}{在真实业务场景下使用具体技术解决具体问题，补充可验证指标，例如采纳率提升、Hit Rate 提升、延迟降低或成本下降。}
    \end{itemize}

    \vspace{0.45em}

    \section{\sectionicon{\faProjectDiagram}{项目经历}}
    \projectheading{项目名称}{项目角色}{2025.10--2026.02}
    \begin{itemize}
        \resumeitem{项目目标}{用一句话说明业务场景、数据规模、目标用户和系统价值。}
        \resumeitem{方案设计}{说明个人负责的技术链路、架构设计、关键取舍和工程实现。}
        \resumeitem{效果评估}{说明验证方式、数据集、指标、Bad Case 回流或上线结果；没有真实指标时保守描述并标注待验证。}
    \end{itemize}

    \vspace{0.35em}
    \projectheading{第二个项目名称}{项目角色}{2025.06--2025.09}
    \begin{itemize}
        \resumeitem{项目目标}{压缩项目背景，突出问题和业务价值。}
        \resumeitem{技术实现}{写清楚使用的技术、负责的模块和解决的问题。}
        \resumeitem{落地结果}{写上线、复用、用户、性能、稳定性或奖项等可信证明。}
    \end{itemize}

    \vspace{0.45em}

    \section{\sectionicon{\faTrophy}{奖项证书}}
    \begin{itemize}
        \resumeitem{证书}{CET-6、职业证书或其他与岗位相关证书}
        \resumeitem{竞赛}{高认可度、强相关的竞赛奖项，按含金量排序}
        \resumeitem{论文/专利}{论文、专利或科研成果，注明状态时保持真实}
    \end{itemize}

    % 如缺少证据，删除自我评价；如保留，每条都必须有证明。
    % \vspace{0.35em}
    % \section{\sectionicon{\faUserCheck}{自我评价}}
    % \begin{itemize}
    %     \resumeitem{AI 应用落地能力}{用真实项目证明需求拆解、方案设计、服务封装、部署和评估能力。}
    % \end{itemize}

\end{document}
